% ID: SEL-D0015-BC-E06
% ESTADO: validada. S016. Revisión matemática, didáctica y visual.
\textbf{a) Probabilidad de buena calidad.}\par\nopagebreak[4]
Las empresas forman una partición con pesos $P(E_1)=0.30$ y $P(E_2)=0.70$. Sea $G$ buena calidad. Por probabilidad total,
\[
\boxed{P(G)=0.30\cdot0.80+0.70\cdot0.60=0.66.}
\]
\textbf{b) Empresa primera y mala calidad.}\par\nopagebreak[4]
En $E_1$ la calidad mala tiene probabilidad $1-0.80-0.05=0.15$. Por la regla del producto,
\[
\boxed{P(E_1\cap\text{mala})=0.30\cdot0.15=0.045.}
\]
\textbf{c) Calidad media condicionada a E2.}\par\nopagebreak[4]
Llamamos $q=P(\text{media}\mid E_2)$. El nuevo dato global permite despejarlo por probabilidad total:
\[
0.19=0.30\cdot0.05+0.70q
\quad\Longrightarrow\quad
\boxed{q=\frac{0.19-0.015}{0.70}=0.25.}
\]
No se reparte por igual el porcentaje desconocido entre media y mala. Como comprobación, en $E_2$ la mala es $1-0.60-0.25=0.15$. La tabla conjunta queda
\[
\begin{array}{c|rrr|r}
&\text{Buena}&\text{Media}&\text{Mala}&\text{Total}\\\hline
E_1&0.24&0.015&0.045&0.30\\
E_2&0.42&0.175&0.105&0.70\\\hline
\text{Total}&0.66&0.19&0.15&1
\end{array}
\]
Esta tabla usa el dato adicional de c; sus marginales y todas sus entradas son coherentes.
